\chapter*{ABSTRAK}
\addcontentsline{toc}{chapter}{ABSTRAK}

\begin{center}

    \textbf{Perancangan Sistem \textit{Business Intelligence} Berbasis \textit{Rule-Based Chatbot} dengan \textit{Time Series Forecasting} dan Teknologi Bahasa Alami untuk Analisis dan Prediksi Kinerja Bisnis} \\
    \vspace{0.4cm}
    oleh \\
    \vspace{0.4cm}
    Thalita Zahra Sutejo\\
    18222023\\
\end{center}


\begin{spacing}{1.0}
Sistem \textit{Business Intelligence} konvensional di lingkungan internal perusahaan umumnya masih mengandalkan pelaporan manual dan memerlukan keahlian teknis dalam bahasa kueri, sehingga menghambat pengguna nonteknis dalam memperoleh \textit{insight} bisnis secara cepat dan efisien. Penelitian ini bertujuan untuk merancang dan mengimplementasikan sistem \textit{Business Intelligence} berbasis \textit{rule-based chatbot} yang mengintegrasikan tiga komponen utama: klasifikasi maksud berbasis pencocokan pola berbobot (\textit{weighted keyword matching}), pembangkitan bahasa alami berbasis templat (\textit{template-based Natural Language Generation}), dan peramalan deret waktu berbasis pembelajaran mendalam (\textit{deep learning}) menggunakan arsitektur GRU (\textit{Gated Recurrent Unit}) dan LSTM (\textit{Long Short-Term Memory}). Sistem dibangun dengan arsitektur tiga lapis yang terdiri dari \textit{frontend} berbasis React.js, \textit{backend} berbasis Node.js/Express yang terhubung ke \textit{Enterprise Data Warehouse} Oracle, serta layanan \textit{machine learning} berbasis FastAPI/PyTorch. Mesin aturan (\textit{rule engine}) mendukung lebih dari 30 aturan analitik yang mencakup lima domain data bisnis, yaitu penjualan, pendapatan, target, proses data pelanggan, dan \textit{churn}. Modul NLG menghasilkan respons naratif dalam bahasa Indonesia dengan akurasi faktual 100\% karena seluruh data numerik berasal langsung dari hasil kueri basis data. Model peramalan deret waktu dievaluasi menggunakan protokol \textit{walk-forward validation} dengan metrik sMAPE, MASE, dan \textit{Skill Score} terhadap \textit{seasonal naive baseline}. Hasil evaluasi menunjukkan bahwa sistem mampu memproses kueri analitik dan menghasilkan respons naratif yang informatif, serta memberikan prediksi indikator bisnis pelanggan kepada pengguna internal perusahaan melalui antarmuka percakapan berbasis web.

Kata kunci: \textit{Business Intelligence}, \textit{Chatbot}, \textit{Rule-Based}, \textit{Natural Language Generation}, \textit{Time Series Forecasting}, GRU, LSTM.
\end{spacing}


